\documentclass[11pt, a4paper,titlepage]{report}
% !BIB program = biber

\input{mypackages}
\makeatletter
\renewcommand\thesection{\arabic{section}}
\renewcommand\thesubsection{\thesection.\arabic{subsection}}
\renewcommand\thesubsubsection{\thesubsection.\arabic{subsubsection}}
\makeatother


\begin{document}
%Personnalisation des listes 
\renewcommand{\labelitemi}{$\bullet$} 
\input{sections/umons-coverpage-test}
\tableofcontents 

\vspace{1em}

\section*{Executive Summary}
This document describes the implementation, validation and evaluation of a
"server congestion resume" extension for QUIC. The extension enables a client
to present to a server the most recently advertised congestion control state
from a previous connection, allowing the server to adopt a warmed congestion
window. The implementation is provided in the modified `quiche` codebase. The
report explains design choices, security considerations, test results and the
cryptographic binding used to protect against forgery and large-scale replay
without server-side storage.

\section{Introduction}
The project implements a QUIC transport parameter (`enable_server_congestion_resume`)
and two new frames: `CC_INDICATION` and `CC_RESUME`. The server periodically
advertises its internal congestion control snapshot using `CC_INDICATION`. The
client stores the most recent advertisement and can send it once at the start
of a subsequent connection using `CC_RESUME`. The server may accept or reject
the proposed state.

This report documents the implementation, how to run the client and server,
security mitigations against forgery and replay, and an initial quantitative
evaluation in a high BDP (bandwidth-delay product) scenario.

\section{Design and Implementation}
The chosen representation for the congestion control state is compact and
portable. The state contains the server wall time, the congestion window in
bytes, the smoothed RTT in nanoseconds, the delivery rate in bits per second,
and a one-byte algorithm tag identifying the congestion control implementation
(Reno, CUBIC or BBRv2). The payload format is a simple binary versioned blob
with a 4-byte magic header `QCCR` and a single version byte.

To provide integrity and authenticity, the state is authenticated using
HMAC-SHA512. The MAC is computed over the concatenation of a obfuscated `epoch`
value and the `cc_state` blob. The implementation supports configuring the
HMAC key; a hard-coded default key exists for testing only. All verification
uses a constant-time comparison. The hash length on the wire is dynamic and is
validated against the received `Hash Length` field, allowing future migration
to different hash sizes without additional server-side state.

To avoid long-lived replay attacks without server storage, the implementation
adds an epoch obfuscation step and derives a per-epoch key from the master key
and the obfuscated epoch. The derived key is used to compute the HMAC, so the
MAC is bound to the epoch value as sent on the wire. This approach preserves
statelessness while strengthening resistance to large-scale replay.

On the server side the implementation emits `CC_INDICATION` frames only when
the congestion snapshot changes significantly (cwnd change >20\%, delivery rate
change >10\%, RTT change >10\% or cwnd doubling) and enforces a minimum
interval of 50 ms between advertisements. On the client side a single stored
state can be injected with `CC_RESUME` during the first 1-RTT packet of the
connection. The server processes only the first `CC_RESUME` and ignores further
ones.

All code changes were done with minimal disturbance to existing modules. The
new code is implemented under the `cc_resume` module, with small integration
points in the `Connection` send and receive paths.

\section{Security and Cryptographic Binding}
The basic HMAC prevents simple forging of `cc_state` blobs. However, a client
or an on-path attacker could replay a previously observed `CC_INDICATION`
into another connection within the validity window. Because the specification
explicitly forbids server-side storage of the advertised state, the binding
must remain stateless. The implementation addresses this by (1) obfuscating
the epoch using a reversible XOR mask derived from the master key, and (2)
deriving a per-epoch key using HMAC-SHA512 over the obfuscated epoch. The
server recomputes the derived key from the received epoch and verifies the
MAC. This keeps verification stateless and ties the MAC to the on-wire epoch.

The design still relies on expiration to limit replay windows; a validity
period of 24 hours is enforced by comparing the embedded wall-clock value
with the current time. This limit provides a clear upper bound for replay
feasibility while preserving compatibility with multi-connection workflows.

\subsection*{Illustrative Code Snippets}
\begin{verbatim}
// Derive per-epoch key from the obfuscated epoch (stateless)
let derived_key = derive_cc_resume_key_from_epoch(master_key, epoch_obf);

// Compute MAC with HMAC-SHA512 over epoch_obf || cc_state
let mac = compute_cc_resume_mac(&derived_key, epoch_obf, &cc_state);
\end{verbatim}

\section{Tests and Results}
All implemented tests are unit-level and cover normal operation and attack
cases. The main tests included are:
\begin{itemize}
  \item `cc_indication_is_stored_and_resumed`: end-to-end flow verification
  \item `cc_resume_invalid_hash_is_ignored`: forged MAC is ignored
  \item `cc_resume_replay_is_ignored_after_first`: replay of a second
    `CC_RESUME` is ignored after the first application
  \item `cc_resume_expired_state_is_ignored`: expired state is ignored
  \item `cc_resume_epoch_tamper_is_ignored`: epoch tampering breaks MAC
\end{itemize}

Figures below show selected test outputs and packet captures.

\begin{figure}[ht]
  \centering
  \includegraphics[width=0.9\textwidth]{test_ccindication&cc_resume.png}
  \caption{End-to-end functional test. The server emits `CC_INDICATION` in
    1-RTT, the client stores the latest state, and the next connection resumes
    using `CC_RESUME`.}
\end{figure}

\begin{figure}[ht]
  \centering
  \includegraphics[width=0.9\textwidth]{test_cc_indication&cc_resume_2.png}
  \caption{Advertisement timeline. The epoch increases monotonically and is
    obfuscated on the wire; the server derives the MAC key from the received
    epoch to validate the resume frame.}
\end{figure}

\begin{figure}[ht]
  \centering
  \includegraphics[width=0.9\textwidth]{tests_robustness.png}
  \caption{Robustness tests. Invalid hashes, expired states, replayed
    `CC_RESUME`, and epoch tampering are all ignored without closing the
    connection.}
\end{figure}

\begin{figure}[ht]
  \centering
  \includegraphics[width=0.9\textwidth]{Extension enabled verif wireshark.png}
  \caption{Wire-level trace. The transport parameter negotiation is visible
    and the `CC_INDICATION` frame appears in a 1-RTT packet with a hash length
    matching the HMAC-SHA512 output.}
\end{figure}

\section{Quantitative Evaluation}
An initial analysis for a high BDP scenario (satellite-like link with 500 ms
RTT and 100 Mbps bandwidth) is provided. With resume enabled, restoring a
several-megabyte congestion window immediately removes the slow-start ramp
that would otherwise consume multiple RTTs. A conservative example in the
report shows a potential order-of-magnitude improvement in utilization during
the first second after connection establishment. The report contains a short
validation plan to measure these gains using traffic shaping and repeated runs.

\section{Integration with QLOG}
For effective debugging and validation the transport parameter `enable_server_congestion_resume`
must be recorded in QLOG as a named field rather than an unknown parameter.
This makes traces readable and avoids misclassification by qlog consumers. In a
subsequent commit the QLOG schema will be extended with a boolean field in the
`ParametersSet` structure and the frame enum will include `CcIndication` and
`CcResume` event variants.

\section{Limitations and Future Work}
The current implementation enforces the core security checks (HMAC verification,
expiration, algorithm matching and single-apply policy) and derives a per-epoch
MAC key from the obfuscated epoch. The following improvements remain:
\begin{itemize}
  \item Add QLOG schema changes so the transport parameter and frames are
    logged as named fields (not unknown parameters).
  \item Implement server-side rate limiting for repeated CC_RESUME attempts
    to reduce DoS risk.
  \item Extend evaluation with an automated test harness and traffic shaping
    to produce the throughput curves used in the quantitative analysis.
\end{itemize}

\printbibliography[heading=bibintoc,title={Bibliography}]
\end{document}
